En un ensayo simulado, la temperatura del proceso es \(R(t)=20+2t\), en grados Celsius, durante \(0\le t\le10\) segundos. Un sensor inicialmente a \(20\,{}^\circ\mathrm C\) tiene capacidad térmica \(C>0\). La potencia térmica que recibe del proceso es proporcional a la diferencia entre la temperatura del proceso y la temperatura del propio sensor, con coeficiente \(H>0\); \(C/H=3\,\mathrm s\). No hay otros intercambios y el sensor no altera la temperatura del proceso.

Elija las cantidades con las que describirá la medición, formule el modelo con su condición inicial y obtenga la lectura del sensor. Se admite una diferencia de, como máximo, \(4\,{}^\circ\mathrm C\) entre la lectura y la temperatura del proceso: determine durante qué parte del ensayo se cumple esa tolerancia. Si se quiere mantener una rampa durante los diez segundos, ¿cuál es la mayor pendiente constante \(a\ge0\) admisible en \(R(t)=20+at\), con el mismo estado inicial? Justifique la decisión a partir del modelo.
